% Part: normal-modal-logic
% Chapter: tableaux
% Section: countermodels

\documentclass[../../../include/open-logic-section]{subfiles}

\begin{document}

\olfileid{nml}{tab}{cou}

\olsection{\usetoken{P}{tableau} থেকে প্রতিমডেল}

\begin{explain}
  সম্পূর্ণতা উপপাদ্যের প্রমাণে শুধু
  $\Entails !A$ থেকে $\Proves !A$ পাওয়া যায় না;
  $\Entails/ !A$ হলে $!A$-র প্রতিমডেল নির্মাণের
  পদ্ধতিও পাওয়া যায়।
  % BN-SRC-385: source omits ! in the non-validity formula.
  $\Log{K}$-এর ক্ষেত্রে এটি একটি \emph{সিদ্ধান্ত-পদ্ধতি}।
  কারণ $\Entails/ !A$ ধরা যাক।
  \olref[cpl]{prop:complete-tableau}-এর প্রমাণ
  সম্পূর্ণ !!{tableau} নির্মাণের উপায় দেয়, এবং
  $\Log{K}$-এর জন্য তা সর্বদাই সসীম ধাপে শেষ হয়।
  এর প্রস্তাবনামূলক বিধিগুলি কেবল কম জটিলতার
  পূর্বসূচিযুক্ত !!{formula}s যোগ করে। ফলে শাখায়
  $\sFmla{S}{!A}[\sigma]$-র জন্য প্রতিটি
  প্রস্তাবনামূলক বিধি একবারই প্রয়োগ করতে হয়।
  নতুন পূর্বসূচি তৈরি হয় কেবল
  \iftag{prvBox}{$\TRule{\False}{\Box}$}{}\iftag{notprvBox,notprvDiamond}{}{
    ও }\iftag{prvDiamond}{$\TRule{\True}{\Diamond}$}{}
  \iftag{notprvBox,notprvDiamond}{বিধিতে}{বিধিগুলিতে};
  এদেরও একবার প্রয়োগেই একটি নতুন পূর্বসূচি তৈরি হয়।
  \iftag{prvBox}{$\TRule{\True}{\Box}$}{}\iftag{notprvBox,notprvDiamond}{}{
    ও }\iftag{prvDiamond}{$\TRule{\False}{\Diamond}$}{}
  \iftag{notprvBox,notprvDiamond}{বিধিটি}{বিধিগুলি}
  একাধিকবার প্রয়োগের প্রয়োজন হতে পারে, তবে প্রতিটি
  পূর্বসূচির জন্য একবার; নতুন পূর্বসূচিও সসীমসংখ্যক।
  অতএব সসীম ধাপের পরে নির্মাণে হয় বন্ধ শাখা,
  নয়তো সম্পূর্ণ শাখা পাওয়া যায়।

  খোলা সম্পূর্ণ শাখাযুক্ত একটি !!{tableau} নির্মিত হলে
  \olref[cpl]{thm:tableau-completeness}-এর প্রমাণ
  মূল পূর্বসূচিযুক্ত !!{formula}s-সেটকে পরিতৃপ্ত করা
  একটি নির্দিষ্ট মডেল দেয়। তাই $\Gamma \Entails !A$
  হলে যেমন বন্ধ !!{tableau} ও $\Gamma \Proves !A$
  পাওয়া যায়, তেমনি পদ্ধতিগত অনুসন্ধানে সম্পূর্ণ খোলা
  শাখা পাওয়া গেলে $\Gamma \Entails/ !A$ জানা যায়
  এবং প্রতিমডেলও নির্মাণ করা যায়।
\end{explain}

\iftag{prvBox}{
\begin{ex}
  আমরা জানি, $\Proves/ \Box(p \lor q) \lif (\Box p \lor \Box
  q)$। একটি ট্যাবলো নির্মাণ শুরু হয় এভাবে:
  \begin{oltableau}
    [\pFmla{\False}{\Box(p \lor q) \lif (\Box p \lor \Box q)}{1},
      just = \TAss, checked
      [\pFmla{\True}{\Box(p \lor q)}{1},
        just = {\TRule{\False}{\lif}[1]},
        [\pFmla{\False}{\Box p \lor \Box q}{1},
          just = {\TRule{\False}{\lif}[1]}, checked
          [\pFmla{\False}{\Box p}{1},
            just = {\TRule{\False}{\lor}[3]}, checked
            [\pFmla{\False}{\Box q}{1},
              just = {\TRule{\False}{\lor}[3]}, checked
              [\pFmla{\False}{p}{1.1},
                just = {\TRule{\False}{\Box}[4]}, checked
                [\pFmla{\False}{q}{1.2},
                  just = {\TRule{\False}{\Box}[5]}, checked
                ]
              ]
            ]
          ]
        ]
      ]
    ]
  \end{oltableau}
  !!{tableau} এখনও শেষ হয়নি। পরের ধাপে টিকচিহ্নবিহীন
  একমাত্র পঙ্‌ক্তিটি, অর্থাৎ $2$ নম্বর পঙ্‌ক্তির
  পূর্বসূচিযুক্ত !!{formula}
  $\sFmla{\True}{\Box(p \lor q)}[1]$ বিবেচনা করি।
  শাখায় ব্যবহৃত প্রত্যেক পূর্বসূচি, অর্থাৎ $1.1$ ও
  $1.2$ উভয়ের জন্য $\TRule{\True}{\Box}$ প্রয়োগ করি:
  \begin{oltableau}
    [\pFmla{\False}{\Box(p \lor q) \lif (\Box p \lor \Box q)}{1},
      just = \TAss, checked
      [\pFmla{\True}{\Box(p \lor q)}{1},
        just = {\TRule{\False}{\lif}[1]},
        [\pFmla{\False}{\Box p \lor \Box q}{1},
          just = {\TRule{\False}{\lif}[1]}, checked
          [\pFmla{\False}{\Box p}{1},
            just = {\TRule{\False}{\lor}[3]}, checked
            [\pFmla{\False}{\Box q}{1},
              just = {\TRule{\False}{\lor}[3]}, checked
              [\pFmla{\False}{p}{1.1},
                just = {\TRule{\False}{\Box}[4]}, checked
                [\pFmla{\False}{q}{1.2},
                  just = {\TRule{\False}{\Box}[5]}, checked
                  [\pFmla{\True}{p \lor q}{1.1},
                    just = {\TRule{\True}{\Box}[2]}
                    [\pFmla{\True}{p \lor q}{1.2},
                      just = {\TRule{\True}{\Box}[2]}
                    ]
                  ]
                ]
              ]
            ]
          ]
        ]
      ]
    ]
  \end{oltableau}
  এখন 2, 8 ও 9 নম্বর পঙ্‌ক্তিতে টিকচিহ্ন নেই।
  কিন্তু নতুন পূর্বসূচি যোগ হয়নি। তাই 8 ও 9
  নম্বর পঙ্‌ক্তিতে $\TRule{\True}{\lor}$ প্রয়োগ করি,
  উৎপন্ন প্রত্যেক খোলা শাখায়:
  \begin{oltableau}
    [\pFmla{\False}{\Box(p \lor q) \lif (\Box p \lor \Box q)}{1},
      just = \TAss, checked
      [\pFmla{\True}{\Box(p \lor q)}{1},
        just = {\TRule{\False}{\lif}[1]}, checked
        [\pFmla{\False}{\Box p \lor \Box q}{1},
          just = {\TRule{\False}{\lif}[1]}, checked
          [\pFmla{\False}{\Box p}{1},
            just = {\TRule{\False}{\lor}[3]}, checked
            [\pFmla{\False}{\Box q}{1},
              just = {\TRule{\False}{\lor}[3]}, checked
              [\pFmla{\False}{p}{1.1},
                just = {\TRule{\False}{\Box}[4]}, checked
                [\pFmla{\False}{q}{1.2},
                  just = {\TRule{\False}{\Box}[5]}, checked
                  [\pFmla{\True}{p \lor q}{1.1},
                    just = {\TRule{\True}{\Box}[2]}, checked
                    [\pFmla{\True}{p \lor q}{1.2},
                      just = {\TRule{\True}{\Box}[2]}, checked
                      [\pFmla{\True}{p}{1.1},
                        just = {\TRule{\True}{\lor}[8]}, checked, close
                      ]
                      [\pFmla{\True}{q}{1.1},
                        just = {\TRule{\True}{\lor}[8]}, checked
                        [\pFmla{\True}{p}{1.2},
                          just = {\TRule{\True}{\lor}[9]}, checked]
                        [\pFmla{\True}{q}{1.2},
                          just = {\TRule{\True}{\lor}[9]}, checked, close]
                      ]
                    ]
                  ]
                ]
              ]
            ]
          ]
        ]
      ]
    ]
  \end{oltableau}
  একটি খোলা শাখা অবশিষ্ট থাকে, এবং সেটি সম্পূর্ণ।
  তার থেকে মডেলটির জগৎ
  $W = \{1, 1.1, 1.2\}$ নিই: এগুলিই খোলা
  শাখার পূর্বসূচি। অভিগম্যতা-সম্পর্ক
  $R = \{\tuple{1, 1.1}, \tuple{1, 1.2}\}$।
  11 নম্বর পঙ্‌ক্তিতে $\sFmla{\True}{p}[1.2]$
  থাকায় $V(p) = \{1.2\}$; আর 10 নম্বর
  পঙ্‌ক্তিতে $\sFmla{\True}{q}[1.1]$ থাকায়
  $V(q) = \{1.1\}$। মডেলটি
  \olref{fig:counter-Box}-এ আঁকা আছে। যাচাই করলে
  দেখা যায়, এটি $\Box(p \lor q) \lif
  (\Box p \lor \Box q)$-র একটি প্রতিমডেল।
  \begin{figure}
  \begin{center}
    \begin{tikzpicture}[modal]
      \node[world] (w1) [label={[align=right]right:\mFalse{p}\\ \mFalse{q}}]{$1$};
      \node[world] (w2) [label={[align=right]right:\mFalse{p}\\ \mTrue{q}},
        above left=of w1]{$1.1$};
      \node[world] (w3) [label={[align=right]right:\mTrue{p}\\ \mFalse{q}},
        above right=of w1] {$1.2$};
      \draw[->] (w1) to (w2);
      \draw[->] (w1) to (w3);
    \end{tikzpicture}
  \end{center}
  \caption{$\Box(p \lor q) \lif (\Box p \lor \Box
    q)$-র একটি প্রতিমডেল।}
  \ollabel{fig:counter-Box}
\end{figure}
\end{ex}
}
{
\begin{ex}
  আমরা জানি, $\Proves/ (\Diamond p \land \Diamond q) \lif \Diamond(p
  \land q)$। একটি !!{tableau} নির্মাণ শুরু হয় এভাবে:
  \begin{oltableau}
    [\pFmla{\False}{(\Diamond p \land \Diamond q) \lif \Diamond(p \land q)}{1},
      just = \TAss, checked
      [\pFmla{\True}{\Diamond p \land \Diamond q}{1},
        just = {\TRule{\False}{\lif}[1]}, checked
        [\pFmla{\False}{\Diamond(p \land q)}{1},
          just = {\TRule{\False}{\lif}[1]}
          [\pFmla{\True}{\Diamond p}{1},
            just = {\TRule{\True}{\land}[2]}, checked
            [\pFmla{\True}{\Diamond q}{1},
              just = {\TRule{\True}{\land}[2]}, checked
              [\pFmla{\True}{p}{1.1},
                just = {\TRule{\True}{\Diamond}[4]}, checked
                [\pFmla{\True}{q}{1.2},
                  just = {\TRule{\True}{\Diamond}[5]}, checked
                ]
              ]
            ]
          ]
        ]
      ]
    ]
  \end{oltableau}
  !!{tableau} এখনও শেষ হয়নি। পরের ধাপে টিকচিহ্নবিহীন
  একমাত্র পঙ্‌ক্তিটি, অর্থাৎ $3$ নম্বর পঙ্‌ক্তির
  পূর্বসূচিযুক্ত !!{formula}
  $\sFmla{\False}{\Diamond(p \land q)}[1]$
  বিবেচনা করি। শাখায় ব্যবহৃত প্রত্যেক পূর্বসূচি,
  অর্থাৎ $1.1$ ও $1.2$ উভয়ের জন্য
  $\TRule{\False}{\Diamond}$ প্রয়োগ করি:
  % BN-SRC-386: line 3 has a false diamond, not a true diamond.
  \begin{oltableau}
    % BN-SRC-387: source reverses the implication at this tableau root.
    [\pFmla{\False}{(\Diamond p \land \Diamond q) \lif \Diamond(p \land q)}{1},
      just = \TAss, checked
      [\pFmla{\True}{\Diamond p \land \Diamond q}{1},
        just = {\TRule{\False}{\lif}[1]}, checked
        [\pFmla{\False}{\Diamond (p \land q)}{1},
          just = {\TRule{\False}{\lif}[1]}
          [\pFmla{\True}{\Diamond p}{1},
            just = {\TRule{\True}{\land}[2]}, checked
            [\pFmla{\True}{\Diamond q}{1},
              just = {\TRule{\True}{\land}[2]}, checked
              [\pFmla{\True}{p}{1.1},
                just = {\TRule{\True}{\Diamond}[4]}, checked
                [\pFmla{\True}{q}{1.2},
                  just = {\TRule{\True}{\Diamond}[5]}, checked
                  [\pFmla{\False}{p \land q}{1.1},
                    just = {\TRule{\False}{\Diamond}[3]}
                    [\pFmla{\False}{p \land q}{1.2},
                      just = {\TRule{\False}{\Diamond}[3]}
                    ]
                  ]
                ]
              ]
            ]
          ]
        ]
      ]
    ]
  \end{oltableau}
  এখন 3, 8 ও 9 নম্বর পঙ্‌ক্তিতে টিকচিহ্ন
  নেই। কিন্তু নতুন পূর্বসূচি যোগ হয়নি।
  তাই 8 ও 9 নম্বর পঙ্‌ক্তিতে
  $\TRule{\False}{\land}$ প্রয়োগ করি, উৎপন্ন
  প্রত্যেক খোলা শাখায়:
  \begin{oltableau}
    [\pFmla{\False}{(\Diamond p \land \Diamond q) \lif \Diamond(p \land q)}{1},
      just = \TAss, checked
      [\pFmla{\True}{\Diamond p \land \Diamond q}{1},
        just = {\TRule{\False}{\lif}[1]}, checked
        [\pFmla{\False}{\Diamond(p \land q)}{1},
          just = {\TRule{\False}{\lif}[1]}
          [\pFmla{\True}{\Diamond p}{1},
            just = {\TRule{\True}{\land}[2]}, checked
            [\pFmla{\True}{\Diamond q}{1},
              just = {\TRule{\True}{\land}[2]}, checked
              [\pFmla{\True}{p}{1.1},
                just = {\TRule{\True}{\Diamond}[4]}, checked
                [\pFmla{\True}{q}{1.2},
                  just = {\TRule{\True}{\Diamond}[5]}, checked
                  [\pFmla{\False}{p \land q}{1.1},
                    just = {\TRule{\False}{\Diamond}[3]}, checked
                    [\pFmla{\False}{p \land q}{1.2},
                      just = {\TRule{\False}{\Diamond}[3]}, checked
                      [\pFmla{\False}{p}{1.1},
                        just = {\TRule{\False}{\land}[8]}, checked, close
                      ]
                      [\pFmla{\False}{q}{1.1},
                        just = {\TRule{\False}{\land}[8]}, checked
                        [\pFmla{\False}{p}{1.2},
                          just = {\TRule{\False}{\land}[9]}, checked]
                        [\pFmla{\False}{q}{1.2},
                          just = {\TRule{\False}{\land}[9]}, checked, close]
                      ]
                    ]
                  ]
                ]
              ]
            ]
          ]
        ]
      ]
    ]
  \end{oltableau}
  একটি খোলা শাখা অবশিষ্ট থাকে, এবং সেটি সম্পূর্ণ।
  তার থেকে মডেলটির জগৎ
  $W = \{1, 1.1, 1.2\}$ নিই: এগুলিই খোলা
  শাখার পূর্বসূচি। অভিগম্যতা-সম্পর্ক
  $R = \{\tuple{1, 1.1}, \tuple{1, 1.2}\}$।
  6 নম্বর পঙ্‌ক্তিতে $\sFmla{\True}{p}[1.1]$
  থাকায় $V(p) = \{1.1\}$; আর 7 নম্বর
  পঙ্‌ক্তিতে $\sFmla{\True}{q}[1.2]$ থাকায়
  $V(q) = \{1.2\}$।
  % BN-SRC-388: source misprints the q witness prefix as 1.1.
  মডেলটি \olref{fig:counter-Diamond}-এ আঁকা আছে।
  যাচাই করলে দেখা যায়, এটি
  $(\Diamond p \land \Diamond q) \lif
  \Diamond (p \land q)$-র একটি প্রতিমডেল।
  \begin{figure}
  \begin{center}
    \begin{tikzpicture}[modal]
      \node[world] (w1) [label={[align=right]right:\mFalse{p}\\ \mFalse{q}}]{$1$};
      \node[world] (w2) [label={[align=right]right:\mTrue{p}\\ \mFalse{q}},
        above left=of w1]{$1.1$};
      \node[world] (w3) [label={[align=right]right:\mFalse{p}\\ \mTrue{q}},
        above right=of w1] {$1.2$};
      \draw[->] (w1) to (w2);
      \draw[->] (w1) to (w3);
    \end{tikzpicture}
  \end{center}
  \caption{$(\Diamond p \land \Diamond q) \lif
    \Diamond (p \land q)$-র একটি প্রতিমডেল।}
  \ollabel{fig:counter-Diamond}
\end{figure}
\end{ex}
}

\end{document}
